\chapter*{Agradecimientos}
\thispagestyle{empty}

En este espacio deseo agradecer a las personas más importantes que hicieron posible este proceso, especialmente a miembros de mi familia, y amigos, que me apoyaron durante todo el desarrollo de esta investigación.

Deseo agradecer a mi mamá Viviana Solís y a mi papá Jorge Johnson, son las personas que más admiro en esta vida, desde que era pequeño ellos decidieron priorizarme a mi sobre todo lo demás, incluso si eso significaba sacrificar su propio bienestar, gracias por ayudarme a impulsar mis sueños. 

A mi mejor amigo y compañero de carrera Jose Antonio Espinoza, que ha estado presente en mi vida desde que estamos en la escuela, me acompañó por todo este proceso universitario, y gracias a él decidí estudiar esta carrera al ver el proyecto del carro de Taller de Programación. 

A mis compañeros de carrera, Max Garro, y Mauricio Calderón, que me acompañaron durante todo este proceso de carrera, tanto siendo compañeros de proyectos, como conviertiendose en buenos amigos. Gracias por compartir todo este proceso conmigo y por todos los momentos juntos.

A mis amigos del colegio y previamente compañeros universitarios, Johel, Eitan, Manrique, Carlos Luis, que estuvieron cuando recién empezamos este proceso universitario. 

Quiero agradecer al profesor Luis Alberto Chavarría, que fue quien impulsó la idea de este proyecto, y me ofreció el apoyo por todo este proyecto universitario, y me ofreció una guía por medio del desarrollo de esta investigación.

Agradecer al Instituto Tecnológico de Costa Rica por proveerme acceso a la educación pública y a todos los profesores que contribuyeron en mi formación tanto académica como profesional, transmitiendome ese amor por la ingeniería. 

Finalmente especial agradecimiento a la Escuela de Ingeniería en Computadores, y al profesor Jeff Schmidt Peralta por haber sido el primer profesor en recibirme en la carrera y proveerme las bases del conocimiento y formación para salir adelante.



\vspace*{1cm}

\scriptAuthor

Cartago, \today

\cleardoublepage

%%% Local Variables: 
%%% mode: latex
%%% TeX-master: "paMain"
%%% End: 
